\documentclass[conference]{IEEEtran}
\IEEEoverridecommandlockouts

\usepackage{cite}
\usepackage{amsmath,amssymb,amsfonts}
\usepackage{algorithmic}
\usepackage{graphicx}
\usepackage{textcomp}
\usepackage{xcolor}
\usepackage[hidelinks]{hyperref}
\def\BibTeX{{\rm B\kern-.05em{\sc i\kern-.025em b}\kern-.08em
    T\kern-.1667em\lower.7ex\hbox{E}\kern-.125emX}}
\usepackage{subcaption}
\begin{document}

\title{House Price Prediction Using Machine Learning: A Study on the Impact of Spatial and Environmental Features}\\

\author{
\IEEEauthorblockN{Yuyang Zhang}
\IEEEauthorblockA{
School of Computer Science \\
University of Nottingham Ningbo China \\
Ningbo, China \\
Student ID: 20514470\\
scyyz26@nottingham.edu.cn
}
}

\begin{abstract}
Predicting house prices is one of the main tasks in real estate research. However, house price data are characterized by non-linearity, spatial dependency, and heavy tails; moreover, no model has been found to perform better than others in all cases. One can use stacking ensembles to combine complementary algorithms. Still, the approach to predicting house prices is inconsistent in the literature, and most studies do not reveal the contribution of each base learner when building an ensemble. In this paper, an end-to-end pipeline is developed on the King County housing data (4,600 observations, 18 variables) to answer the following questions: does stacking yield better results compared to the best-performing individual model, and which base learners have the most influence on performance improvements? To begin with, six base learners are evaluated against the same preprocessing procedure: Ridge, Random Forest, Extra Trees, GBR, XGBoost, and Multi-Layer Perceptron (MLP). Then, three stacking schemes are built on the same base learners using out-of-fold predictions and either Ridge or LassoCV as the meta-model. The best stacking model yields R² = 0.7429, which is only marginally better than the best base model, XGBoost (R² = 0.7375). Leave-one-out ablation shows that Random Forest (+0.0030) and XGBoost (+0.0011) are mostly responsible for performance improvements; other base learners make no substantial contribution, and Ridge even harms ensemble performance. Therefore, increasing the number of base learners might not be the right way to improve ensemble performance since it is mostly determined by several strong learners and spatial features coded as the target variable mean.
\end{abstract}

\begin{IEEEkeywords}
house price prediction, stacking ensemble, ablation study, target encoding, gradient boosting
\end{IEEEkeywords}
\section{Introduction}
House prices are a crucial concern for buyers, sellers, lenders, and policymakers \cite{KOK2014136}, and they are notoriously difficult to model. Each property constitutes an almost unique combination of structures, location, and condition, where the relationship between each factor and price is nonlinear and imbalanced across the market \cite{35659ca8-6be8-375a-94eb-6f497d2fc083}. Methods based on additive linear relationships, like hedonic regression and other parametric models, struggle to account for the interactions and give inconsistent prices when transaction numbers are sparse.  \cite{KOK2014136}.

Machine learning techniques relax these assumptions and are now ubiquitous methods for property valuation. Nevertheless, none of these methods performs consistently across all datasets and markets. Ridge, linear models in general, are easy to interpret but incapable of modeling non-linearities \cite{limsombunchao2004house}. On the other hand, tree-based ensembles perform well in modeling interactions but suffer from over-fitting and interpretability issues \cite{fan2006determinants,RICOJUAN2021114590} . Deep networks are highly flexible but architecture choices and tuning affect results \ cite {10984370,JIT2774}. This problem motivates stacking, where a meta-model combines the outputs of several heterogeneous learners so as to keep the best of each one. Yet, stacking is not extensively used in real estate price prediction due to inconsistent advantages reported so far\cite{land13091436}. Most of the related literature uses single learners and simple combinations of learners, such as bagging and boosting algorithms \cite{RICOJUAN2021114590}. The few studies that use stacking show some gains but treat the ensemble as a black-box approach with little discussion of what each of the base learners contributes to the prediction task \cite{Bjrgve2025ComparingHV} . Without this understanding, it is uncertain if the improvements generalize or a simpler ensemble would work just as well.

This study tried to fill the gap using the King County housing dataset (4,600 transactions, 18 features, May 2014 to May 2015). In particular, it combines structural attributes (number of bedrooms, bathrooms, and square feet of living space), environmental characteristics (waterfront and view), and high cardinality location features (city and state/zip code), which allows various learners to take advantage of each piece of information differently. Contributions are:
\begin{itemize}
    \item A controlled comparison of six baseline regressors and three stacking strategies with identical preprocessing.
    \item A leave-one-out experiment analyzing how important each base learner is to the stacked model through its contribution to the final prediction, aided by meta-learner coefficients.
    \item  Identify that stacking produces the most value in Random Forest and XGBoost, and that target means of location features are dominant in feature importance. 
\end{itemize}

\section{Literature Review}
Data-driven approaches to estimating the value of houses can be classified into three broad groups: parametric baselines, single-learner machine learning, and ensemble models. 

For parametric baseline models, Hedonic regression models assume a linear relationship between the price and characteristics of a property  \cite{KOK2014136}. They are straightforward and cost-effective to apply, but due to their linear assumption, they are susceptible to non-linearity and spatial heterogeneity in the data \cite{SELIM20092843,limsombunchao2004house}.

As machine learning gained traction in this area, a range of single-learner models has been applied to housing valuation. Tree-based learners — decision trees \cite{fan2006determinants}, random forests \cite{10.1023/A:1010933404324}, and gradient boosting \cite{hjort2022house,10.1214/aos/1013203451} — are the most widely used, because they do not require feature scaling and capture feature interactions without explicit transformation of the target. Neural networks have outperformed linear baselines in some comparisons \cite{SELIM20092843} but are prone to overfitting and offer limited interpretability \cite{limsombunchao2004house}. Support vector regression with a non-linear kernel has also been applied, but is sensitive to the choice of CC C and the ϵ\textbackslash{}epsilon ϵ-insensitive loss \textbackslash{}cite\{CHERKASSKY2004113\}. The recurring conclusion across these studies is that any of these learners can be made to work well on a given dataset with careful tuning, but no single family is consistently best.

This is the gap that ensemble methods are designed to close. Bagging and boosting reduce the variance of individual learners by aggregating many of them \cite{10.1023/A:1010933404324,hjort2022house}. Stacked generalisation goes further by training a meta-learner on the predictions of heterogeneous base models \cite{WOLPERT1992241}, and is in principle better placed than a single ensemble family to exploit complementary errors. In the housing-valuation literature, however, stacking remains relatively rare. The studies that do use it tend to focus on overall accuracy without controlled comparisons against individually tuned baselines \cite{land13091436,10937620}, and rarely decompose the ensemble to show which base learners actually contribute. Bjrgve et al. \cite{Bjrgve2025ComparingHV} note this explicitly, arguing that the practical case for stacking in housing valuation is still weak because gains are often modest and the underlying mechanism is opaque.


\section{Methodology}
The pipeline has six stages: exploratory analysis, cleaning, feature engineering, feature selection, baseline modelling, and stacking. Each stage is described below, with leakage-aware design choices flagged where they apply.
\subsection{Exploratory Analysis}
The dataset is loaded from \texttt{data.csv} and inspected for missing values, dtype consistency, and target skew. No imputation is needed. The price distribution is right-skewed with a long upper tail, so outlier handling is required before training error-sensitive metrics such as RMSE. Continuous predictors are visualised against price (regression plots), categorical predictors via boxplots, and pairwise relationships via a correlation heatmap.

\subsection{Outlier Removal}
Price observations beyond the 99th percentile are dropped (46 observations, 1.0\% of the data). Such a decision removes a few extremely valuable and luxury houses that would dominate model fitting and RMSE objective.

\subsection{Feature Engineering}
\texttt{date} column is used for \textit{house\_age}, \textit{is\_renovated}, and \textit{years\_since\_renovation} extraction, thus defining property age. \textit{basement\_ratio} captures the relative size of the basement in the house, whereas \textit{has\_basement} is a binary indicator; a small constant ensures no division by zero. \textit{bath\_bed\_ratio} contains information on the number of bathrooms and bedrooms. \textit{month} is converted to the number of season and then encoded via one-hot encoding.

\subsection{Encoding and Train–Test Split}
After the application of feature engineering techniques, raw identifiers (\textit{street}, \textit{city}, and \textit{statezip}) are dropped. Data split into training ($80\%$, 3,643 instances) and test ($20\%$, 911 instances) sets is done randomly but with a fixed random seed. To avoid the dimensional growth and at the same time preserve location information, categorical \textit{city} and \textit{statezip} columns are reintroduced through target encoding: categories become replaced by the average price of train set members, whereas unknown categories of the test set are assigned to the average price across all categories from the train set. Average prices are calculated using only training data to eliminate leakage. 

\subsection{Feature Selection and Scaling}
The feature selection step combines automatic and manual techniques. Features are selected automatically based on the criteria of correlation with price: $|r| > 0.05$. In addition, a manual set of domain-relevant features is saved regardless of their univariate correlation. Features with high mutual correlation, $|r|>0.85$, among other retained features are pruned. Two exceptions from this rule are \textit{sqft\_living}, \textit{bathrooms}, and the target-encoded location columns. As a result of this operation, only 14 features are left (dropped is \textit{basement\_ratio}). Remaining features are standardized based on training data only.

\subsection{Baseline Models}
Six baseline models belong to three groups: regularized linear, bagging, and boosting ensembles. Linear ridge regression acts as the representative of the first group. Bagging ensembles include Random Forest and Extra Trees. GBR and XGBoost stand for boosting ensemble techniques. Finally, the neural network regressor (MLP) provides non-linear representation capabilities. All models share the same scaled feature set and are evaluated with $R^{2}$, RMSE, and MAE.

\subsection{Stacking Ensemble}
Out-of-fold (OOF) predictions are generated for each base learner using 5-fold cross-validation: each fold's held-out predictions form the meta-features for that fold, and test predictions are averaged across folds. Three stacking configurations are then trained on these meta-features. Configuration (\textbf{V1} )  uses only the OOF predictions, utilizes a Ridge meta-learner. In the second configuration (\textbf{V2}) the meta-features are concatenated with the original, scaled features before feeding into the Ridge meta-learner. Lastly, the third configuration (\textbf{V3}) substitutes LassoCV for the meta-learner, thus introducing automatic regularization and feature selection.

\subsection{Evaluation}
All models are assessed using the held-out 20\% test split. To reach the second research goal, a leave-one-out ablation on the top performing stacking approach is performed, whereby the impact of each base learner on the test  $R^{2}$ is measured by removing one base learner at a time. Additionally, meta-learner coefficients for V1 and V3 are provided for two reasons. First, they show how the predictions from different base learners are combined and indicate any redundancy in the models through correlation of the coefficients. Second, for configurations what including Ridge meta-leanrer, V1 outperforms than V2(See details in Results). Random Forest feature importances and three diagnostic plots (predicted-versus-actual, residuals against predictions, residual histogram) are used to inspect bias and heteroscedasticity.

\section{Experimental Results}
The results will be presented in the same order as the pipeline above: observations from EDA, followed by an evaluation of all models, an ablation study, meta-learner coefficients, and diagnostics.

\subsection{EDA and Feature Engineering Insights}
\begin{figure}[htbp]
\centering
\includegraphics[width=\linewidth]{images/output.png}
\caption{Price distribution and outlier analysis. (a) Right-skewed original distribution. (b) Boxplot before outlier removal. (c) After dropping the top 1\% (46 rows), the bulk of the distribution is far more compact.}
\label{fig:price_outlier}
\end{figure}

The price distribution is an indication of the necessity of dealing with outliers. The price distribution shown in Fig. ~\ref{fig:price_outlier}a is extremely skewed to the right, where most data points are clustered on the left side while a number of points are scattered far out to the right. This is presented graphically using a box plot in Fig. ~\ref{fig:price_outlier}b. The price distribution in Fig. ~\ref{fig:price_outlier}c is compressed after removing the top 1\% of the data points.

\begin{figure}[htbp]
\centering
\includegraphics[width=\linewidth]{images/eda_relationships.png}
\caption{Key predictors against price. (a) Living area is positively correlated with price. (b) Waterfront properties command a clear premium. (c) Higher view ratings track higher prices monotonically.}
\label{fig:eda_relationships}
\end{figure}

Fig.\ref{fig:eda_relationships} depicts the three main factors that guided our feature engineering efforts. Panel (a) shows a plot between \textit{sqft\_living} and price, showing that there is an increasing trend, as well as increasing variation for higher values of this feature, implying non-linear nature. Panel (b) shows that properties that have water fronts form a group with higher prices.

\begin{figure}[htbp]
\centering
\includegraphics[width=\linewidth]{images/feature_engineering.png}
\caption{Engineered features against price. House age, renovation status, and season individually carry weak signal but contribute non-additively to ensemble models.}
\label{fig:feature_engineering}
\end{figure}

These constructed features, as shown in Fig. \ref{fig:feature_engineering}, have very little predictive power by themselves. The correlation shown in Fig. \ref{fig:feature_engineering}a is a weak negative one between the age of the house and the price. Fig. \ref{fig:feature_engineering}b shows that there is only a slight difference between the prices of houses that have been renovated and those that have not been. Fig. \ref{fig:feature_engineering}c shows that there is not much variation depending on the season.

\begin{figure}[htbp]
\centering
\includegraphics[width=\linewidth]{images/city_price.png}
\caption{Average price by city. The distribution is long-tailed: a few cities account for the high end of the market, justifying target encoding rather than one-hot.}
\label{fig:city_price}
\end{figure}

City effects (Fig.~\ref{fig:city_price}) are sizable and heavy-tailed, as just a few cities are responsible for the higher end. With numerous cities included in the dataset, one-hot encoding would be redundant since it increases the dimensionality of the feature space without adding any value.

\begin{figure}[htbp]
\centering
\includegraphics[width=\linewidth]{images/feature_correlation_with_price.png}
\caption{Absolute correlation of each retained feature with price. \textit{zip\_te}, \textit{sqft\_living}, and \textit{city\_te} are the three strongest individual predictors.}
\label{fig:feature_corr}
\end{figure}

The ranking of correlation shown in Figure \ref{fig:feature_corr} gives the inputs that have contributed the most to the signal. The top three predictors include zip\_te ($r$=0.668), sqft\_living ($r$=0.640), and city\_te ($r$=0.532). The secondary set of predictors comprises bathrooms ($r$=0.489), bedrooms ($r$=0.327), and view ($r$=0.321). While the engineered features rank relatively low, they are kept according to the rules outlined above. The importance of both the location variables encoded by target encoding stands out as a common theme for the rest of the analysis.

\begin{figure}[htbp]
\centering
\includegraphics[width=\linewidth]{images/corr_comparison.png}
\caption{Correlation matrix before and after feature selection. (a) Strong multicollinearity among size-related features(left). (b) After dropping \textit{basement\_ratio}, redundant pairs are removed and relationships with the target are preserved(right).}
\label{fig:corr_comparison}
\end{figure}

The before-and-after comparison (see Fig.~\ref{fig:corr_comparison}) illustrates the effect of the multicollinearity reduction phase: only the predictor {basement\_ratio} is eliminated, leaving a 14-predictor model in which all the meaningful correlations are retained. Therefore, linear models are robust in further analysis.

\subsection{Model Comparison}

\begin{table}[htbp]
\centering
\caption{Test-set performance of all nine models, sorted by $R^{2}$.}
\label{tab:model_performance}
\begin{tabular}{lccc}
\hline
Model & $R^{2}$ & RMSE & MAE \\
\hline
Stacking V1 & 0.7429 & 147405 & 88005 \\
Stacking V2 & 0.7428 & 147437 & 88232 \\
Stacking V3 & 0.7422 & 147608 & 88009 \\
XGBoost & 0.7375 & 148950 & 88548 \\
GBR & 0.7298 & 151094 & 89896 \\
Random Forest & 0.7276 & 151724 & 88525 \\
MLP & 0.7046 & 158004 & 99704 \\
Extra Trees & 0.7026 & 158540 & 94566 \\
Ridge & 0.7002 & 159160 & 100678 \\
\hline
\end{tabular}
\end{table}

\begin{figure}[htbp]
\centering
\includegraphics[width=\linewidth]{images/model_comparison.png}
\caption{Test $R^{2}$ for all nine models. Stacking variants narrowly beat the best individual model.}
\label{fig:model_comparison}
\end{figure}

The comparison between the models is shown in Table \ref{tab:model_performance} and Figure \ref{fig:model_comparison}. From the individual models, XGBoost performs the best ($R^{2}$ = 0.7375) compared to GBR and Random Forest. The models using boosting and bagging algorithms rank high among the individual models, while the linear and neural baseline models are ranked relatively low. The result aligns with the finding from the exploratory data analysis (EDA) showing that the most important features have non-linear relationships with each other. All stacking configurations beat XGBoost with a margin of superiority in terms of $R^{2}$ and RMSE. Among the three configurations, V1 shows the highest $R^{2}$ score of 0.7429. The difference in $R^{2}$ and RMSE scores between the highest performing model (V1) and XGBoost is about 0.005 and 1,500, respectively. Although the margins are small, the improvements show consistency in all three performance measures and all three stacking configurations, suggesting that the differences are not accidental. Comparing V1 to the other two stacking configurations, we find that Configuration V1 performs better with an advantage in terms of $R^{2}$ of less than 0.001.

\subsection{Ablation and Meta-Learner Coefficients}

\begin{figure}[htbp]
\centering
\includegraphics[width=0.85\linewidth]{images/ablation.png}
\caption{Leave-one-out ablation on Stacking V1. A positive $\Delta R^{2}$ means removing the model hurts performance.}
\label{fig:ablation}
\end{figure}

\begin{table}[htbp]
\centering
\caption{Leave-one-out ablation on Stacking V1.}
\label{tab:ablation}
\begin{tabular}{lcc}
\hline
Removed model & Resulting $R^{2}$ & $\Delta R^{2}$ \\
\hline
Random Forest & 0.7399 & $+0.0030$ \\
XGBoost       & 0.7418 & $+0.0011$ \\
Extra Trees   & 0.7424 & $+0.0005$ \\
GBR           & 0.7427 & $+0.0002$ \\
MLP           & 0.7428 & $+0.0000$ \\
Ridge         & 0.7433 & $-0.0004$ \\
\hline
None (full)   & 0.7429 & $0$ \\
\hline
\end{tabular}
\end{table}

\begin{table}[htbp]
\centering
\caption{Meta-learner coefficients on the OOF predictions. V1 uses Ridge; V3 uses LassoCV (the Extra Trees coefficient was set to exactly zero by Lasso).}
\label{tab:meta_coefs}
\begin{tabular}{lcc}
\hline
Base learner   & V1 (Ridge) & V3 (Lasso) \\
\hline
Random Forest  & $+0.5630$  & $+0.4503$  \\
XGBoost        & $+0.2610$  & $+0.2612$  \\
Ridge          & $+0.2415$  & $+0.2277$  \\
GBR            & $+0.0802$  & $+0.0900$  \\
MLP            & $+0.0172$  & $+0.0093$  \\
Extra Trees    & $-0.1306$  & $0$ (dropped) \\
\hline
\end{tabular}
\end{table}

As can be seen from the ablation results in Fig. \ref{fig:ablation} and summarized in Table \ref{tab:ablation}, the hierarchy is obvious: the main player here is Random Forest, whose elimination causes a decrease in R-squared value of $\Delta R^{2}=0.0030$, about three times greater than the one caused by elimination of XGBoost ($\Delta R^{2}=0.0011$). GBR, MLP, and Extra Trees' contributions do not exceed $0.0005$, and removing Ridge leads to the marginal improvement of $0.0004$. This hierarchy, as expected, is supported by the V1 coefficients presented in Table \ref{tab:meta_coefs}. The largest coefficients correspond to the most important base models: Random Forest has a positive coefficient of $+0.563$, and XGBoost – $+0.261$. Ridge also gets a significant coefficient: $+0.242$. GBR and MLP get smaller values: $+0.080$ and $+0.017$, respectively, whereas Extra Trees gets negative $-0.131$ value. This does not mean that, isolated, Extra Trees model would perform worse. It uses the same algorithm as Random Forest, but the split criteria are different. As a consequence, the correlation between Extra Trees' and Random Forest's predictions is high, so that the Ridge learner uses Extra Trees in order to remove the overlap and leave only disagreement, which is the goal of stacking technique.

Another evidence supporting the idea of redundancy of Extra Trees comes from the meta-Lasso in V3 setting. When allowed to drop predictors, it gives zero weight to Extra Trees, while the other five base models obtain nearly equal coefficients to those obtained previously in V1. Since the new model gets almost equal R-squared value ($R^{2}=0.7422$) compared to V1's result ($\Delta R^{2}=0.0008$), it shows that Extra Trees is redundant given that Random Forest is included in the ensemble. Thus, the conclusion we can draw based on the above results and analysis is that stacking mainly leverages Random Forest and XGBoost, whereas other learners have marginal effects, and Extra Trees turns out to be fully redundant in the presence of Random Forest model.

\subsection{Diagnostics and Without Preprocessing Training}

\begin{figure}[htbp]
\centering
\includegraphics[width=\linewidth]{images/model_diagnostics_stacking_v3.png}
\caption{Diagnostics for Stacking V3: predicted vs actual, residuals vs predicted, and residual distribution.}
\label{fig:diagnostics}
\end{figure}

The results shown in Figure \ref{fig:diagnostics} confirm what is presented in the table. As can be seen from Figure \ref{fig:diagnostics}a, the predictions are tightly clustered around the line of identity. Figure \ref{fig:diagnostics}b shows that the residuals are centered around zero but their variance increases with increasing predicted values, which implies some degree of heteroscedasticity since the model has lower precision when predicting the prices of more costly properties, a result that is not unexpected since there were few observations left after filtering out those beyond the 99th percentile threshold.

\begin{figure}[htbp]
\centering
\includegraphics[width=0.8\linewidth]{images/model_performance_before_after_preprocessing.png}
\caption{$R^{2}$ for Random Forest and XGBoost before and after applying the preprocessing pipeline.}
\label{fig:preprocess}
\end{figure}

Figure~\ref{fig:preprocess} shows the effect of pre-processing. Without pre-processing, both Random Forest and XGBoost produce an R² on test set of about 0.03. After the pre-processing step is applied to the data, the Random Forest produces a score of 0.723 and the XGBoost model 0.733. The difference in the results is considerably greater than the best to worst result difference shown in Table~\ref{tab:model_performance}.

\section{Discussion}
The outcomes show that ensemble models, and especially those based on trees, better fit the structure of housing data compared to linear or neural models. This can be seen from the ranking provided in Table~\ref{tab:model_performance}: XGBoost, Gradient Boosting Regressor (GBR), and Random Forest are among the top three, while Ridge regression and the MLP are at the bottom of the list. This outcome is coherent with the EDA results: the relationship between \textit{sqft\_living} and price is non-linear; waterfront and view variables are responsible for price steps; and target-encoded location features represent categorical information about price quite well. Linear models do not capture this non-linearity due to the additive assumption, which is not satisfied for housing data, while the MLP requires a larger sample and/or more expressive features to outperform the boosting baselines~\cite{hjort2022house,10.1214/aos/1013203451}.

Furthermore, stacking provides a further improvement, even though it is minor. The three stacking variants exceed XGBoost by about $0.005$ in $R^{2}$, an increase consistently observed across all three stacking configurations. This phenomenon is clarified by the ablation analysis. Among the six learners, Random Forest and XGBoost provide the most valuable input to the ensemble; Extra Trees mainly decorrelates Random Forest, while the contribution of the remaining three learners is minimal. As a result, the meta-learner does not integrate six different perspectives on the problem but rather reweights two tree-based perspectives. This outcome is consistent with the findings of Bj{\o}rgve et al.~\cite{Bjrgve2025ComparingHV}, according to which the gains obtained through stacking in housing valuation are moderate when the stronngest base learners is not the right direction to optimize as base -learner diversity is what really counts. 

A second observation concerns how the contribution of each base learner varies depending on whether it works in isolation or stacked with other learners. For example, XGBoost outperforms Random Forest by $0.010$ in $R^{2}$ as a standalone model, but when included in the stacking ensemble, Random Forest contributes nearly three times as much as XGBoost (Table~\ref{tab:ablation}). This indicates that the contribution made by a base learner to the ensemble is not determined by its standalone performance but by the complementary information it carries relative to the other learners. The OOF predictions of XGBoost apparently overlap with those of GBR and Extra Trees, which use related algorithmic mechanisms. On the contrary, Random Forest contributes a pattern difficult to replicate by the meta-learner using other learners' predictions. Such information cannot be captured by performance metrics; that is why leave-one-out ablation helps identify it.

The last two comments apply to all model comparisons performed during this research project. First, the analysis of the preprocessing variant shows that the choice of data preparation strategy has a greater impact on model test performance than model selection itself: from $R^{2}\approx 0.03$ to $R^{2}>0.72$ upon removing outliers and engineering features. This outcome implies that, contrary to the conventional practice, preprocessing should not be considered an automatically performed preliminary step but rather regarded as a modelling decision. Second, the dominance of the target-encoded location features in the correlation ranking , \textit{zip\_te} (r=0.668, r=0.668, r=0.668) and \textit{city\_te} (r=0.532,r=0.532, r=0.532) relating to \textit{sqft\_living} (r=0.640r=0.640
r=0.640) , confirms the importance of spatial information in housing valuation. It can be explained by the well-known problem of ignoring spatial heterogeneity in hedonic models \cite{SELIM20092843}.

\section{Conclusion and Future Work}
This paper put a stacking ensemble for house price prediction through a controlled comparison and a leave-In this study, we analyze stacking for house price predictions through a controlled experiment on the King County Housing dataset. Our optimal stacking setup reached $R^{2}$ = 0.7429, barely beating XGBoost ($R^{2}$ = 0.7375). In addition, the ablation analysis shows that the contribution to the result from our stacking setup is virtually entirely made by Random Forest and XGBoost. With Lasso meta-learner, it also seems clear that this result would be nearly identical if the last four algorithms were left out. This means that there are just two learners worth including in the stack and adding other learners does not contribute much. This suggests, practically, that optimizing stack performance should involve choosing good base learners, not increasing the stack size.

Two additional results follow from this analysis. The first is that preprocessing has a larger impact than model selection on this dataset. Tree-based learners move from $R^{2}\approx 0.03$ on raw data to $R^{2}>0.72$ once outliers are removed and engineered features are added. This shows that the signal connecting features to price exists in the dataset, but only becomes accessible to the model after preprocessing. The second result is that the target-encoded location features (\textit{zip\_te}, \textit{city\_te}) rank among the strongest predictors, on a par with \textit{sqft\_living}. Geography therefore accounts for a substantial share of the price variance, and target encoding offers a low-cost way to capture this signal without the dimensionality cost of one-hot encoding.

This study has three main limitations. First, the model performs poorly on luxury houses due to insufficient representation by 99th- percentile filtering and lack of suitable features. Second, we use purely structural tabular features to make predictions. Including some other signals like economic factors, schools, geographic coordinates etc. could have provided further gains. Third, adding more learners did not lead to significant performance boosts as the gains seem to come mostly from a few good learners. These limitations point to three directions for future research. The first one is to train the model with custom loss functions designed to handle the right part of the distribution and/or weight high-cost cases. This could mitigate the heteroscedastic issue\cite{hjort2022house}. The second option is to experiment with different types of geographic encoding like learned embeddings or coordinate features instead of using target-encoding. Finally, instead of doing ablation studies after meta-learner fitting, we need a principled method for selecting base learners based on diversity.














% This study adopts a structured end-to-end data modelling pipeline for house price prediction. The pipeline includes exploratory data analysis, data preprocessing, feature engineering, feature selection, baseline model construction, stacking-based ensemble learning, and model evaluation. Each stage is designed to improve predictive performance, reduce noise, and minimise the risk of data leakage.\\

% The analysis begins with exploratory data analysis (EDA) to gain an overall understanding of the dataset structure and the relationships between variables. The dataset is first loaded from \texttt{data.csv} and examined through summary statistics, data types, and sample records, providing an initial overview of its characteristics. To assess data quality, missing values are analysed by computing both their counts and proportions across all features. As no significant missingness is observed, no imputation is required.Attention is then directed towards the target variable (\textit{price}). Its distribution is visualised using a histogram with kernel density estimation, and skewness is calculated to quantify asymmetry. This step helps identify potential issues such as imbalance or extreme values. Building on this, the relationships between key predictors and house price are explored. Regression plots are used to examine continuous variables such as \textit{sqft\_living}, while boxplots are employed to analyse the influence of categorical variables, including \textit{waterfront} and \textit{view}. Finally, a correlation heatmap is constructed for numerical variables, allowing important predictors to be identified and potential multicollinearity to be detected. Together, these steps provide a comprehensive understanding of the dataset and inform subsequent modelling decisions.\\

% To reduce the impact of extreme values, outliers are removed by excluding observations above the 99th percentile of \textit{price}. This improves model stability and is particularly important for error-sensitive metrics such as RMSE.\\

% Following data cleaning, feature engineering is conducted to enrich the representation of the dataset and better capture underlying patterns in house prices. The \textit{date} variable is first converted into a datetime format, enabling the extraction of temporal attributes such as \textit{year} and \textit{month}. These temporal features are then used to derive additional variables, including \textit{house\_age}, \textit{is\_renovated}, and \textit{years\_since\_renovation}, which reflect the lifecycle characteristics of properties. In addition to temporal transformations, several structural features are introduced to describe the internal composition of houses. For instance, \textit{basement\_ratio} captures the proportion of basement space relative to total living area, while \textit{has\_basement} provides a binary indication of its presence. Similarly, \textit{bath\_bed\_ratio} represents the balance between bathrooms and bedrooms, offering a more informative measure than absolute counts. A small constant is incorporated in ratio calculations to ensure numerical stability and avoid division-by-zero issues. Finally, to account for potential seasonal effects in the housing market, the month variable is mapped into four categories, including winter, spring, summer, and autumn, forming a new \textit{season} feature. This allows the model to capture temporal variations in demand that may influence property prices.\\

% After feature engineering, irrelevant and redundant variables are removed. Raw location and identifier variables such as \textit{street}, \textit{city}, and \textit{statezip} are excluded at this stage, as they are either high-cardinality or replaced by engineered representations. The categorical \textit{season} feature is encoded using one-hot encoding, and boolean variables are converted into numeric form to ensure compatibility with machine learning models. The dataset is then split into training and testing sets using an 80/20 ratio with a fixed random seed.\\

% To preserve spatial information without introducing high dimensionality, target encoding is applied to \textit{statezip} and \textit{city}. The mean house price for each category is computed using only the training data and then mapped to both training and test sets. Unseen categories in the test set are assigned the global mean price. This approach captures location-specific effects while preventing data leakage.\\

% Feature selection is performed using a hybrid strategy combining statistical filtering and domain knowledge. Features with an absolute correlation greater than 0.05 with the target variable are retained. In addition, several domain-relevant features are manually preserved to ensure that important predictors are not discarded. To address multicollinearity, pairwise correlations among selected features are analysed, and highly correlated variables (correlation above 0.85) are removed using an upper-triangular filtering approach. However, key variables such as \textit{sqft\_living}, \textit{bathrooms}, and the target-encoded features are protected to maintain model interpretability.\\

% Before modelling, feature scaling is applied using standardisation. The scaler is fitted on the training data and applied to both training and test sets to avoid leakage. This step is particularly important for models sensitive to feature magnitude, such as Ridge regression and neural networks.\\

% A different set of baseline models is then constructed to capture different learning behaviours. Ridge regression is used as a regularised linear baseline. Random Forest and Extra Trees represent bagging-based ensemble methods capable of modelling nonlinear interactions. Gradient Boosting and XGBoost are included as boosting-based approaches that iteratively reduce residual errors. In addition, an MLP regressor is used to capture complex nonlinear patterns through neural network representation learning. All models are trained on the scaled features and evaluated using $R^2$, RMSE, and MAE to provide complementary performance perspectives.\\

% To further enhance performance, a stacking ensemble framework is implemented. Five-fold cross-validation is used to generate out-of-fold (OOF) predictions for each base model. For each fold, the model is trained on a subset of the data and used to predict the held-out fold, producing unbiased meta-features for training. Test predictions are generated in each fold and averaged to ensure stability.\\

% Three stacking configurations are explored. The first uses only OOF predictions as inputs to a Ridge meta-learner. The second extends this by introducing a passthrough mechanism, where original scaled features are concatenated with OOF predictions, allowing the meta-learner to utilise both sources of information. The third replaces the Ridge meta-learner with LassoCV, enabling automatic regularisation and feature selection within the ensemble. This design improves flexibility and allows the model to adaptively select informative predictors.\\

% Model performance is evaluated on the held-out test set using $R^2$, RMSE, and MAE, and results are compared across all baseline and ensemble models. To analyse the contribution of each base model, an ablation study is conducted by removing one model at a time from the stacking ensemble and measuring the resulting performance drop. This provides insight into the relative importance of each component. In addition, feature importance is assessed using a Random Forest model, and diagnostic plots including predicted-versus-actual, residual scatter, and residual distribution are used to evaluate model behaviour and identify potential biases.\\

% Overall, the proposed methodology integrates systematic preprocessing, domain-informed feature engineering, leakage-aware encoding, hybrid feature selection, diverse baseline models, and stacking-based ensemble learning. These design choices collectively improve model robustness, generalisation, and interpretability.\\

% \section{Experimental Results and Analysis}
% This section presents the experimental process and results of the proposed modelling pipeline. It begins with exploratory data analysis and feature engineering insights to understand data characteristics and guide feature design. Subsequently, multiple regression models are implemented and compared, including both baseline methods and stacking-based ensembles, with performance evaluated using standard metrics.\\

% \subsection{Exploratory Analysis and Feature Engineering Insights}
% \begin{figure}[htbp]
% \centering
% \includegraphics[width=\linewidth]{images/output.png}
% \caption{Distribution and outlier analysis of house prices. (a) The original price distribution is highly right-skewed, indicating the presence of extreme values. (b) Before outlier removal, a large number of extreme high-price observations are visible. (c) After removing the top 1\% of prices, the distribution becomes more concentrated and stable.}
% \label{fig:price_outlier}
% \end{figure}
% The experimental process begins with an exploratory analysis of the dataset to understand the distribution of the target variable and the relationships between key features. As shown in Fig.~\ref{fig:price_outlier}, the distribution of house prices is highly right-skewed, with a dense concentration of lower-valued properties and a long tail of extreme high prices. This is further confirmed by the boxplot, where numerous outliers extend far beyond the upper quartile. Such extreme values can disproportionately affect model training and lead to unstable predictions, particularly for error-sensitive metrics such as RMSE. \\

% To address this issue, a percentile-based filtering strategy is applied by removing the top 1\% of price values. After preprocessing, the distribution becomes more compact with reduced variance, indicating improved data stability for subsequent modelling.
% \begin{figure}[htbp]
% \centering
% \includegraphics[width=\linewidth]{images/eda_relationships.png}
% \caption{Relationships between key features and house price. (a) Living area shows a positive correlation with price. (b) Properties with waterfront access have significantly higher prices. (c) Price increases with higher view ratings, indicating a monotonic trend.}
% \label{fig:eda_relationships}
% \end{figure}
% Following this, the relationships between key predictors and house price are examined in Fig.~\ref{fig:eda_relationships}. A clear positive relationship is observed between \textit{sqft\_living} and price, although the dispersion suggests non-linear effects. In addition, properties with waterfront access and higher view ratings consistently achieve higher prices, indicating that both structural and location-related features are important determinants of house value. These observations motivate the inclusion of such variables in the modelling stage.
% \begin{figure}[htbp]
% \centering
% \includegraphics[width=\linewidth]{images/feature_engineering.png}
% \caption{Impact of engineered features on house prices. (a) House age shows a weak negative relationship with price. (b) Renovated properties exhibit slightly higher average prices. (c) Seasonal variation indicates marginal differences across time periods.}
% \label{fig:feature_engineering}
% \end{figure}
% To further enhance the representational capacity of the dataset, several engineered features are introduced and evaluated in Fig.~\ref{fig:feature_engineering}. Features such as \textit{house\_age}, \textit{is\_renovated}, and seasonal indicators show relatively weak individual relationships with price. However, these variables are retained as they may provide complementary information or capture interaction effects when used in more complex models.
% \begin{figure}[htbp]
% \centering
% \includegraphics[width=\linewidth]{images/city_price.png}
% \caption{Average house price across different cities. Significant variation in average prices is observed, indicating strong location-dependent effects.}
% \label{fig:city_price}
% \end{figure}
% Spatial variation is further analysed in Fig.~\ref{fig:city_price}, where substantial differences in average house prices are observed across cities. The ranked distribution reveals a long-tailed pattern, indicating that a small number of locations contribute disproportionately to high property values. This confirms that location is a key determinant of price. Given the high cardinality of spatial variables, direct one-hot encoding would be inefficient; therefore, target encoding is adopted to preserve location information while controlling feature dimensionality.
% \begin{figure}[htbp]
% \centering
% \includegraphics[width=\linewidth]{images/feature_correlation_with_price.png}
% \caption{Absolute correlation between features and house price. Strong correlations are observed for \textit{sqft\_living}, \textit{zip\_te}, and \textit{city\_te}, while several engineered features exhibit moderate or weak relationships.}
% \label{fig:feature_corr}
% \end{figure}
% To quantify feature importance, the absolute correlation between each feature and the target variable is presented in Fig.~\ref{fig:feature_corr}. The results show that \textit{sqft\_living} and the target-encoded spatial features (\textit{zip\_te} and \textit{city\_te}) have the strongest associations with price, followed by variables such as \textit{bathrooms}, \textit{bedrooms}, and \textit{view}. In contrast, several engineered features exhibit weaker correlations but are retained due to their potential contribution in non-linear or ensemble models.
% \begin{figure}[htbp]
% \centering
% \includegraphics[width=\linewidth]{images/corr_comparison.png}
% \caption{Comparison of correlation structures before and after feature selection. (a) The initial correlation matrix shows strong multicollinearity among several features (e.g., \textit{sqft\_living} and \textit{sqft\_above}). (b) After feature selection, redundant variables are removed, resulting in a more balanced and less correlated feature space.}
% \label{fig:corr_comparison}
% \end{figure}
% Finally, the effectiveness of feature selection is evaluated by comparing the correlation structures before and after refinement, as shown in Fig.~\ref{fig:corr_comparison}. The initial feature set exhibits strong multicollinearity, particularly among size-related variables (e.g., \textit{sqft\_living} and \textit{sqft\_above}), with correlations reaching up to 0.88. After applying correlation-based filtering, redundant variables are removed, resulting in a more balanced and less correlated feature space. Importantly, key relationships with the target variable are preserved. This reduction in multicollinearity improves model stability and interpretability, particularly for linear models.

% \subsection{Model Performance and Comparative Evaluation}

% Following the exploratory analysis and feature engineering stage, the refined feature set is used to train and evaluate multiple regression models. These include a linear model (Ridge), several tree-based ensemble methods (Random Forest, Extra Trees, Gradient Boosting, and XGBoost), a neural network model (MLP), and stacking-based ensembles. This section presents a comparative evaluation of these approaches, followed by an analysis of prediction behaviour and the impact of preprocessing.

% \begin{table}[htbp]
% \centering
% \caption{Performance comparison of models}
% \label{tab:model_performance}
% \begin{tabular}{lccc}
% \hline
% Model & R$^2$ & RMSE & MAE \\
% \hline
% Stacking V1 & 0.7467 & 146316 & 88025 \\
% Stacking V3 & 0.7462 & 146440 & 87915 \\
% Stacking V2 & 0.7439 & 147113 & 88716 \\
% XGBoost & 0.7417 & 147741 & 88430 \\
% GBR & 0.7366 & 149207 & 89691 \\
% Random Forest & 0.7295 & 151188 & 88370 \\
% Extra Trees & 0.7053 & 157818 & 93898 \\
% MLP & 0.7018 & 158732 & 100606 \\
% Ridge & 0.7003 & 159141 & 100720 \\
% \hline
% \end{tabular}
% \end{table}

% \begin{figure}[htbp]
% \centering
% \includegraphics[width=\linewidth]{images/model_comparison.png}
% \caption{Comparison of model performance in terms of R$^2$. Stacking-based ensembles consistently outperform individual models.}
% \label{fig:model_comparison}
% \end{figure}

% The overall performance results are summarised in Table~\ref{tab:model_performance} and visualised in Fig.~\ref{fig:model_comparison}. Among the individual models, XGBoost achieves the highest $R^2$ score (0.7417), followed by Gradient Boosting and Random Forest. This confirms the effectiveness of tree-based ensemble methods in modelling nonlinear relationships and complex feature interactions.\\

% Stacking models further improve performance beyond all individual approaches. In particular, Stacking V1 achieves the best overall performance ($R^2 = 0.7467$), while Stacking V3 performs comparably. Stacking V2 shows slightly lower performance, likely due to differences in the meta-learning configuration. Although the improvement over the best individual model is moderate, the consistent gains across all metrics demonstrate that stacking effectively combines complementary model behaviours. This suggests that the ensemble is better able to capture both linear and nonlinear patterns than any single model.

% \begin{figure}[htbp]
% \centering
% \includegraphics[width=\linewidth]{images/model_diagnostics_stacking_v3.png}
% \caption{Prediction performance and residual analysis of a representative stacking model (Stacking V3).}
% \label{fig:diagnostics}
% \end{figure}

% To further assess model behaviour, Fig.~\ref{fig:diagnostics} presents diagnostic plots for a representative stacking model. The predicted-versus-actual plot shows that most observations lie close to the diagonal line, indicating strong predictive accuracy. \\

% The residual plot does not exhibit a clear systematic pattern, suggesting that the model captures the main structure of the data. However, the dispersion of residuals increases for higher predicted values, indicating mild heteroscedasticity. The residual distribution is centred around zero, confirming that the model does not exhibit significant bias. Overall, the results indicate that the stacking model provides reliable predictions, with slightly reduced accuracy at extreme price levels.

% \begin{figure}[htbp]
% \centering
% \includegraphics[width=0.8\linewidth]{images/model_performance_before_after_preprocessing.png}
% \caption{Impact of preprocessing on model performance.}
% \label{fig:preprocess}
% \end{figure}

% The importance of preprocessing is highlighted in Fig.~\ref{fig:preprocess}. A substantial improvement in performance is observed after applying data cleaning and feature engineering, with $R^2$ increasing from near zero to above 0.73 for both Random Forest and XGBoost. This indicates that model performance is primarily constrained by data quality rather than model complexity, and that appropriate preprocessing is essential for effective learning.

% \begin{figure}[htbp]
% \centering
% \includegraphics[width=0.85\linewidth]{images/ablation.png}
% \caption{Ablation study showing the contribution of each base model to the stacking ensemble. Positive $R^2$ drop indicates performance degradation when the model is removed.}
% \label{fig:ablation}
% \end{figure}

% To analyse the contribution of individual base learners, an ablation study is conducted as shown in Fig.~\ref{fig:ablation}. Each model is removed from the ensemble in turn, and the resulting change in performance is measured.\\

% The results indicate that removing XGBoost leads to the largest drop in performance ($\Delta R^2 = 0.0021$), highlighting its dominant contribution within the ensemble. Random Forest and Extra Trees also contribute positively, although their impact is comparatively smaller.\\

% In contrast, removing Ridge and Gradient Boosting results in a slight improvement in performance, suggesting that these models may introduce redundant or overlapping information when combined with stronger learners. The contribution of MLP is negligible, as its removal produces almost no change in performance.\\

% Overall, the ablation analysis demonstrates that the effectiveness of the stacking model is primarily driven by a small subset of strong base learners, particularly XGBoost. This highlights the importance of carefully selecting complementary models rather than simply increasing ensemble size.\\

% In summary, the experimental results demonstrate that model performance is strongly influenced by both data quality and model design. Effective preprocessing and feature engineering significantly improve predictive capability, while ensemble methods further enhance performance by integrating complementary information from multiple models. Among all approaches, stacking achieves the best results, with its effectiveness largely driven by a small number of strong base learners, particularly XGBoost. These findings highlight the importance of combining robust data preparation with carefully selected models to achieve reliable and accurate house price prediction.

% \section{Discussion}

% The experimental results provide several important insights into the effectiveness of different modelling approaches for house price prediction. Rather than focusing solely on performance metrics, this section interprets the observed patterns and examines the underlying reasons for model behaviour.\\

% First, the strong performance of tree-based ensemble methods, particularly XGBoost and Gradient Boosting, highlights the importance of capturing nonlinear relationships and feature interactions in housing data. As observed in the exploratory analysis, key variables such as \textit{sqft\_living}, \textit{view}, and location-based features exhibit complex and non-linear effects on price. Linear models such as Ridge regression, which assume additive and linear relationships, are therefore inherently limited in their ability to model such patterns. This explains their comparatively lower performance. These findings are consistent with prior studies in house price prediction, which report that boosting-based models often outperform linear approaches due to their flexibility in modelling complex data structures.\\

% Second, the improvement achieved by stacking models suggests that different algorithms capture complementary aspects of the data. While tree-based models are effective at modelling nonlinearities, linear models can still contribute by stabilising predictions, and neural networks may capture subtle patterns not explicitly represented in tree splits. However, the relatively modest gain from stacking indicates that the strongest individual models already capture most of the underlying signal. This is further supported by the ablation analysis, where the removal of XGBoost leads to the largest performance drop, confirming its dominant role within the ensemble. In contrast, removing Ridge or Gradient Boosting slightly improves performance, suggesting redundancy or overlapping information among certain base learners.\\

% Third, the diagnostic analysis reveals that although the stacking model achieves strong overall accuracy, its performance is not uniform across all price ranges. The presence of heteroscedasticity, with larger residual variance at higher price levels, indicates that expensive properties remain more difficult to predict. This can be attributed to the limited number of high-value samples and the greater variability in luxury housing characteristics, which are not fully captured by the available features. Similar challenges have been reported in previous research, where models tend to perform worse on extreme values due to data imbalance and increased noise.\\

% Furthermore, the significant improvement observed after preprocessing demonstrates that data quality plays a more critical role than model complexity. The removal of outliers and the introduction of engineered features substantially improve model performance, transforming the dataset into a more learnable representation. This supports the widely recognised principle in machine learning that effective feature engineering and data cleaning are essential prerequisites for high-performing models.\\

% Finally, the results also highlight the importance of feature representation, particularly for spatial variables. The strong contribution of target-encoded location features confirms that geographic information is a key determinant of house prices. Compared to traditional encoding methods, target encoding provides a compact and informative representation, enabling models to effectively utilise high-cardinality categorical variables without introducing excessive dimensionality.\\

% Overall, the findings suggest that successful house price prediction depends on three key factors: high-quality data preprocessing, appropriate feature representation, and the use of flexible models capable of capturing nonlinear relationships. While ensemble methods such as stacking offer additional improvements, their effectiveness relies on the inclusion of strong and diverse base learners rather than simply increasing model complexity.


% \section{Conclusion and Future Work}

% This study presents a structured end-to-end modelling pipeline for house price prediction, integrating data preprocessing, feature engineering, feature selection, and multiple regression models, including stacking-based ensembles. The experimental results show that model performance is strongly influenced by both data quality and model design.\\

% Among the evaluated approaches, tree-based ensemble models demonstrate superior performance, confirming their ability to capture nonlinear relationships and feature interactions within the data. The stacking framework further improves predictive accuracy by combining complementary strengths of multiple models, although the overall gain is moderate. This indicates that while ensemble learning is effective, its performance largely depends on the inclusion of strong base learners rather than increasing model complexity alone.\\

% The results also highlight the critical importance of preprocessing and feature representation. Techniques such as outlier removal, target encoding, and feature engineering significantly improve model performance, suggesting that well-prepared data is essential for effective learning. In particular, spatial features encoded through target encoding are shown to be key contributors, reinforcing the importance of location in house price prediction.\\

% Despite these strengths, several limitations remain. First, the model exhibits reduced accuracy for high-value properties, as indicated by increased residual variance. This is likely due to data imbalance and the limited representation of extreme cases. Second, the current approach relies solely on structured tabular features and does not incorporate external data sources such as economic indicators or geographic information, which may further enhance predictive performance. Third, the stacking framework, while effective, shows diminishing returns, suggesting potential redundancy among certain base models.\\

% Future work can address these limitations in several directions. Incorporating additional data sources, such as temporal economic trends or spatial coordinates, may improve model generalisation. More advanced feature engineering techniques, including interaction features or learned embeddings, could further enhance representation. In addition, exploring alternative ensemble strategies, such as weighted stacking or model selection based on diversity, may lead to more efficient integration of base learners. Finally, specialised modelling approaches for high-value properties could be developed to better handle heteroscedasticity and improve prediction accuracy in extreme cases.\\

% Overall, this study demonstrates that effective house price prediction requires a combination of high-quality data preprocessing, appropriate feature representation, and robust modelling techniques, with ensemble methods providing additional but controlled performance gains.

\bibliographystyle{IEEEtran}
\bibliography{references}
\vspace{12pt}

\end{document}
